\chapter*{Lời Mở Đầu}
\addcontentsline{toc}{chapter}{Lời Mở Đầu}
\markboth{Lời Mở Đầu}{Lời Mở Đầu}

\begin{truyen}{Lời Mở Đầu}{Opening Words}
\chuhoa{N}{hững chiếc bẫy nguy hiểm nhất không phải do người khác giăng — mà là những chiếc bẫy ta tự bước vào.}
\textit{(The most dangerous traps are not set by others --- they are the ones we step into ourselves.)}

Bẫy của sự tự mãn. Bẫy của cái tôi quá lớn. Bẫy của việc so sánh liên tục. Bẫy của việc trì hoãn mãi không dứt. Quyển mười một này mô tả những bẫy phổ biến nhất mà người trẻ thường tự bước vào — đôi khi không hề hay biết.
\textit{(The trap of complacency. The trap of an inflated ego. The trap of constant comparison. The trap of endless procrastination. Volume eleven describes the most common traps young people walk into willingly --- sometimes without even realizing.)}

Hiểu tên của cái bẫy là cách đầu tiên để thoát khỏi nó.
\textit{(Knowing the name of the trap is the first step to escaping it.)}

Không ai miễn nhiễm — nhưng ai đọc sách này với tâm thế thật thà sẽ thấy mình trong ít nhất một câu chuyện.
\textit{(No one is immune --- but anyone who reads this book with honesty will find themselves in at least one story.)}
\end{truyen}

\begin{baihoc}
Bẫy tự bước vào đau hơn bẫy người khác đặt --- vì nó đi kèm với sự xấu hổ.
\textit{(Traps we walk into ourselves hurt more than those others set --- because they come with shame.)}
\end{baihoc}

\begin{ghinhoanh}
\textbf{Short English to remember:}

Self-deception is the hardest trap to escape.

Awareness is the first exit from every trap.
\end{ghinhoanh}

\begin{tuhocnhanh}
1. Em nhận ra mình đang ở trong chiếc bẫy nào trong quyển này không? \textit{Do you recognize yourself in any trap described in this volume?}

2. Điều gì khiến ta khó thoát khỏi những bẫy ta tự tạo ra? \textit{What makes it hard to escape traps we create for ourselves?}

3. Ai hoặc điều gì đã từng giúp em thoát ra khỏi một bẫy như vậy? \textit{Who or what has helped you escape such a trap before?}
\end{tuhocnhanh}

\ngancach
